\chapter{Esercizi sui fondamenti}

Questa parte raccoglie gli esercizi svolti, organizzati per livello. Molti riprendono gli esempi visti nei capitoli teorici
\textbf{con numeri diversi}, così da potersi esercitare davvero; altri sono le esercitazioni di laboratorio proposte a lezione,
con lo svolgimento completo. Il consiglio è sempre lo stesso: \textbf{provare prima da soli}, e solo dopo leggere la soluzione.

% ══════════════════════════════════════════════════════════════
\section{Ritardi e tempi di trasmissione}
% ══════════════════════════════════════════════════════════════

\info{Richiamo: i ritardi di un pacchetto}{
    Un pacchetto che attraversa un nodo accumula quattro ritardi:
    $$d_{\text{nodo}} = d_{\text{elab}} + d_{\text{coda}} + d_{\text{trasm}} + d_{\text{prop}}$$
    \begin{itemize}
      \item $d_{\text{elab}}$: elaborazione (controllo degli errori, scelta dell'uscita), tipicamente microsecondi;
      \item $d_{\text{coda}}$: attesa in coda, dipende dal traffico;
      \item $d_{\text{trasm}} = L/R$: tempo per ``spingere'' sul collegamento tutti gli $L$ bit al rate $R$. Dipende dalla
            \textbf{dimensione del pacchetto} e dal \textbf{rate}, \emph{non} dalla distanza;
      \item $d_{\text{prop}} = d/s$: tempo che il segnale impiega a percorrere la distanza $d$ alla velocità $s$
            ($2$--$3 \cdot 10^8$ m/s). Dipende dalla \textbf{distanza}, \emph{non} dal pacchetto.
    \end{itemize}
}

\begin{esercizio}{Trasmissione contro propagazione}
Un pacchetto di $L = 1000$ byte viaggia su un collegamento lungo $d = 2500$ km, con rate $R = 10$ Mb/s e velocità di
propagazione $s = 2{,}5 \cdot 10^8$ m/s.
\begin{enumerate}[label=(\alph*)]
    \item Calcolare il tempo di trasmissione e il tempo di propagazione. Quale dei due domina?
    \item Ripetere il calcolo con $R = 1$ Gb/s. Che cosa cambia?
    \item Quanto è ``lungo'', in metri, un singolo bit sul cavo nei due casi?
    \item Un collegamento via satellite geostazionario (quota $36\,000$ km, $s = 3\cdot 10^8$ m/s) collega due stazioni a
          terra. Quanto vale il ritardo di propagazione tra le due stazioni?
\end{enumerate}
\end{esercizio}

\svolgimento
\textbf{(a)} Il pacchetto è di $1000 \cdot 8 = 8000$ bit.
$$d_{\text{trasm}} = \frac{L}{R} = \frac{8000}{10^7} = 0{,}8 \text{ ms} \qquad
  d_{\text{prop}} = \frac{d}{s} = \frac{2{,}5 \cdot 10^6 \text{ m}}{2{,}5 \cdot 10^8 \text{ m/s}} = 10 \text{ ms}$$
Domina la propagazione (12,5 volte più grande); il totale, senza code ed elaborazione, è $10{,}8$ ms.

\textbf{(b)} Con $R = 10^9$ b/s: $d_{\text{trasm}} = 8000/10^9 = 8\ \mu$s, mentre $d_{\text{prop}}$ resta $10$ ms. Il totale scende
solo a $10{,}008$ ms: aumentare la banda \textbf{non accorcia la propagazione}.

\textbf{(c)} In un secondo escono $R$ bit che si distribuiscono su $s$ metri, quindi ogni bit occupa $s/R$ metri:
\begin{itemize}
    \item a 10 Mb/s: $2{,}5\cdot 10^8 / 10^7 = 25$ m per bit (il pacchetto intero è lungo $200$ km);
    \item a 1 Gb/s: $2{,}5\cdot 10^8 / 10^9 = 25$ cm per bit.
\end{itemize}

\textbf{(d)} Il segnale sale e scende: $d = 2 \cdot 36\,000$ km $= 7{,}2\cdot 10^7$ m, quindi
$d_{\text{prop}} = 7{,}2\cdot 10^7 / 3\cdot 10^8 = 0{,}24$ s. Un RTT via satellite supera così il mezzo secondo, ed è per questo
che le telefonate satellitari hanno quel fastidioso ritardo.

\warning{L'analogia dell'autostrada}{
    Pensate a una colonna di 10 auto (i bit) che attraversa un casello (il router) e poi percorre 100 km di autostrada. Il
    tempo per far passare \emph{tutte} le auto dal casello è il ritardo di \textbf{trasmissione}; il tempo che ciascuna
    impiega a percorrere l'autostrada è il ritardo di \textbf{propagazione}. Aprire più caselli (più banda) non rende
    l'autostrada più corta.
}

\begin{esercizio}{Messaggio intero o pacchetti?}
Un messaggio di $8 \cdot 10^6$ bit va dalla sorgente alla destinazione attraverso due router, cioè su $N = 3$ collegamenti
da $R = 2$ Mb/s. Si trascurino propagazione, code ed elaborazione.
\begin{enumerate}[label=(\alph*)]
    \item Quanto tempo serve se il messaggio viene inviato \textbf{intero} (commutazione di messaggio, store-and-forward)?
    \item Quanto tempo serve se viene diviso in $P = 800$ pacchetti da $10\,000$ bit ciascuno? Quando arriva il primo pacchetto?
    \item Ricavare la formula generale e spiegare da dove viene il guadagno.
\end{enumerate}
\end{esercizio}

\svolgimento
\textbf{(a)} Ogni nodo deve ricevere tutto il messaggio prima di inoltrarlo. Ogni collegamento richiede
$8\cdot 10^6 / 2\cdot 10^6 = 4$ s, quindi il totale è $3 \cdot 4 = 12$ s.

\textbf{(b)} Ogni pacchetto richiede $L/R = 10^4 / 2\cdot 10^6 = 5$ ms per collegamento.
\begin{itemize}
    \item Il \textbf{primo} pacchetto arriva dopo $3 \cdot 5 = 15$ ms.
    \item La sorgente finisce di trasmettere l'\textbf{ultimo} pacchetto dopo $800 \cdot 5$ ms $= 4$ s; da lì servono altri 2
          collegamenti, cioè $10$ ms. Il messaggio è completo dopo $4{,}01$ s.
\end{itemize}

\textbf{(c)} La figura mostra il caso di $P = 3$ pacchetti su $N = 3$ collegamenti, con il tempo diviso in intervalli da $L/R$.

\begin{figure}[H]
\centering
\begin{tikzpicture}[x=1.25cm, y=0.8cm, every node/.style={font=\scriptsize\sffamily}]
  \foreach \r/\n in {0/{collegamento 1}, 1/{collegamento 2}, 2/{collegamento 3}} \node[anchor=east] at (0,-\r-0.5) {\n};
  \foreach \t in {0,...,5} \draw[black!25] (\t,0) -- (\t,-3);
  \foreach \t in {1,...,5} \node at (\t-0.5,0.3) {$t_{\t}$};
  \foreach \r in {0,...,3} \draw[black!25] (0,-\r) -- (5,-\r);
  \foreach \r in {0,1,2} {
    \foreach \p/\c in {1/netblue!55, 2/netgreen!55, 3/netorange!60} {
      \pgfmathtruncatemacro{\s}{\r+\p-1}
      \fill[\c] (\s+0.05,-\r-0.1) rectangle (\s+0.95,-\r-0.9);
      \node at (\s+0.5,-\r-0.5) {p\p};
    }
  }
  \draw[decorate, decoration={brace, amplitude=5pt, mirror}] (0,-3.15) -- (5,-3.15) node[midway, below=6pt] {$(P + N - 1)\cdot L/R = (3+3-1)\cdot L/R = 5\cdot L/R$};
\end{tikzpicture}
\caption{Con i pacchetti i collegamenti lavorano \emph{in parallelo}, come una catena di montaggio.}
\end{figure}

Il primo pacchetto impiega $N$ intervalli, e ciascuno dei $P-1$ successivi arriva un intervallo dopo il precedente:
$$T = (P + N - 1)\cdot\frac{L}{R} = (800 + 3 - 1) \cdot 5 \text{ ms} = 4{,}01 \text{ s}$$
Con il messaggio intero, invece, $T = N \cdot P \cdot L/R$: i collegamenti lavorano \emph{uno alla volta}. Dividere in pacchetti
permette ai router di lavorare contemporaneamente; inoltre un errore costringe a ritrasmettere solo un pacchetto, e più
comunicazioni possono alternarsi sullo stesso collegamento. Il prezzo sono gli header di ogni pacchetto e la necessità di
riordinarli a destinazione.

\begin{esercizio}{Throughput e collo di bottiglia}
Un server invia un file di 4 MB a un client attraverso tre collegamenti in serie: $R_1 = 500$ kb/s, $R_2 = 2$ Mb/s,
$R_3 = 1$ Mb/s.
\begin{enumerate}[label=(\alph*)]
    \item Qual è il throughput del trasferimento, in assenza di altro traffico?
    \item Quanto dura, approssimativamente, il trasferimento?
    \item Il collegamento centrale viene condiviso in parti uguali da 10 trasferimenti contemporanei. Come cambiano le risposte?
\end{enumerate}
\end{esercizio}

\svolgimento
\begin{center}
\begin{tikzpicture}[every node/.style={font=\scriptsize\sffamily}]
  \node (s) at (0,0) {\icon[0.7cm]{server}}; \node[below=-1pt of s] {server};
  \node (r1) at (3,0) {\icon[0.8cm]{router}};
  \node (r2) at (6,0) {\icon[0.8cm]{router}};
  \node (c) at (9,0) {\icon[0.8cm]{laptop}}; \node[below=-1pt of c] {client};
  \draw[wired, line width=1pt] (s) -- node[above] {$R_1 = 500$ kb/s} (r1);
  \draw[wired, line width=3pt] (r1) -- node[above] {$R_2 = 2$ Mb/s} (r2);
  \draw[wired, line width=2pt] (r2) -- node[above] {$R_3 = 1$ Mb/s} (c);
  \node[text=netred] at (1.5,-0.45) {collo di bottiglia};
\end{tikzpicture}
\end{center}

\textbf{(a)} Il throughput è limitato dal collegamento più lento: $\min\{R_1, R_2, R_3\} = 500$ kb/s.

\textbf{(b)} $4$ MB $= 32 \cdot 10^6$ bit, quindi $T = 32\cdot 10^6 / 5 \cdot 10^5 = 64$ s.

\textbf{(c)} Ogni trasferimento ottiene $R_2/10 = 200$ kb/s sul collegamento centrale, che diventa il nuovo collo di bottiglia:
throughput $= \min\{500, 200, 1000\} = 200$ kb/s, e il trasferimento dura $32\cdot 10^6 / 2\cdot 10^5 = 160$ s. Il collo di
bottiglia non è sempre il collegamento con il rate più basso: dipende anche da \textbf{quanto è condiviso}.

% ══════════════════════════════════════════════════════════════
\section{Commutazione di circuito e di pacchetto}
% ══════════════════════════════════════════════════════════════

\begin{esercizio}{Quanti utenti su un collegamento}
Un collegamento da 1 Mb/s è condiviso da utenti che, quando sono attivi, trasmettono a 100 kb/s; ogni utente è attivo in
media il 10\% del tempo, in modo indipendente dagli altri.
\begin{enumerate}[label=(\alph*)]
    \item Con la commutazione di circuito, quanti utenti si possono servire?
    \item Con la commutazione di pacchetto e 35 utenti, qual è la probabilità che esattamente $n$ utenti siano attivi?
    \item Qual è la probabilità che più di 10 utenti siano attivi contemporaneamente? Commentare.
\end{enumerate}
\end{esercizio}

\svolgimento
\textbf{(a)} Ogni circuito riserva 100 kb/s anche quando l'utente tace: $1000/100 = 10$ utenti.

\textbf{(b)} Ogni utente è attivo con probabilità $p = 0{,}1$, indipendentemente dagli altri: il numero di utenti attivi segue
una \textbf{distribuzione binomiale}
$$P(n \text{ attivi}) = \binom{35}{n}\, p^n (1-p)^{35-n}$$
Per esempio $P(0) \approx 0{,}025$, $P(3) \approx 0{,}225$ (il valore più probabile), $P(10) \approx 0{,}0013$.

\textbf{(c)} Il collegamento va in sovraccarico solo se ci sono almeno 11 utenti attivi:
$$P(n > 10) = 1 - \sum_{n=0}^{10} \binom{35}{n}\, 0{,}1^n\, 0{,}9^{35-n} \approx 1 - 0{,}99958 = 0{,}00042$$
Circa \textbf{4 volte su diecimila}. La commutazione di pacchetto serve quindi 3,5 volte gli utenti del circuito, con prestazioni
quasi identiche: è il guadagno della \textbf{multiplazione statistica}, che sfrutta il fatto che il traffico dati è
\emph{a raffiche} (\emph{bursty}). Il prezzo è che, in quei rari momenti, i pacchetti si accodano e possono andare persi: non c'è
più una garanzia.

\begin{esercizio}{Un file su un circuito TDM}
Quanto tempo serve per inviare un file di $640\,000$ bit dall'host A all'host B su una rete a commutazione di circuito, se tutti
i collegamenti sono da $1{,}536$ Mb/s, ognuno usa TDM con 24 slot, e servono 500 ms per stabilire il circuito? Quanto
servirebbe se il collegamento fosse tutto per noi?
\end{esercizio}

\svolgimento
Con TDM a 24 slot, a ogni circuito spetta $1{,}536 / 24 = 64$ kb/s. Il trasferimento dura
$640\,000 / 64\,000 = 10$ s, a cui si sommano i $0{,}5$ s di instaurazione: $\mathbf{10{,}5}$ \textbf{s}.

Se potessimo usare l'intero collegamento, la trasmissione durerebbe $640\,000 / 1\,536\,000 \approx 0{,}42$ s (per
collegamento). Il circuito ci \textbf{garantisce} i 64 kb/s, ma ci impedisce di usare la capacità lasciata libera dagli altri.

% ══════════════════════════════════════════════════════════════
\section{Topologie}
% ══════════════════════════════════════════════════════════════

\begin{esercizio}{Contare collegamenti e porte}
Una rete ha $n = 8$ nodi.
\begin{enumerate}[label=(\alph*)]
    \item Quanti collegamenti servono con topologia a maglia completa, a stella, ad anello e a bus? Quante porte deve avere ogni
          nodo nella maglia completa?
    \item Si aggiunge un nono nodo: quanti collegamenti in più servono in ciascun caso?
    \item Che cosa succede, in ciascuna topologia, se si guasta un singolo collegamento? E se si guasta il nodo centrale della stella?
\end{enumerate}
\end{esercizio}

\svolgimento
\textbf{(a)} e \textbf{(b)}:
\begin{table}[H]
\centering
\begin{tabular}{lccc}
\toprule
\textbf{Topologia} & \textbf{Collegamenti con 8 nodi} & \textbf{Con 9 nodi} & \textbf{In più} \\
\midrule
Maglia completa & $\frac{8\cdot 7}{2} = 28$ & $\frac{9 \cdot 8}{2} = 36$ & 8 \\
Stella & 8 (più il nodo centrale) & 9 & 1 \\
Anello & 8 & 9 & 1 (se ne toglie uno e se ne aggiungono due) \\
Bus & 1 dorsale + 8 derivazioni & 1 + 9 & 1 derivazione \\
\bottomrule
\end{tabular}
\caption{Collegamenti richiesti dalle diverse topologie.}
\end{table}
Nella maglia completa ogni nodo è collegato a tutti gli altri: servono $n - 1 = 7$ porte per nodo. Aggiungere un nodo costa
$n$ collegamenti nuovi e una porta in più su \emph{ogni} nodo esistente: è la crescita quadratica che rende la maglia completa
impraticabile per reti grandi.

\textbf{(c)} Guasto di un collegamento:
\begin{itemize}
    \item \textbf{maglia completa}: nessun effetto sulla connettività, esistono altri cammini;
    \item \textbf{stella}: resta isolato solo il nodo di quel collegamento;
    \item \textbf{anello}: se l'anello è unidirezionale la rete si interrompe; se è bidirezionale (o doppio) i messaggi
          possono girare nell'altro verso;
    \item \textbf{bus}: se si guasta la dorsale la rete si divide in due parti (e il segnale riflesso può disturbare tutti).
\end{itemize}
Se si guasta il \textbf{centro della stella} si ferma tutta la rete: è il suo \emph{single point of failure}.

\begin{esercizio}{Perché esistono gli IXP}
Sei ISP regionali vogliono scambiarsi traffico direttamente, a coppie. Quanti collegamenti di peering servono? E se si
incontrano tutti in un IXP? Ripetere con 50 ISP.
\end{esercizio}

\svolgimento
Collegare tutti con tutti è una maglia completa: $\frac{6 \cdot 5}{2} = 15$ collegamenti. Con un IXP ogni ISP porta un solo
collegamento al punto di scambio: $6$ collegamenti. Con 50 ISP: $\frac{50 \cdot 49}{2} = 1225$ contro $50$. L'IXP trasforma una
maglia in una \textbf{stella}, ed è il motivo per cui questi punti neutrali di scambio si sono diffusi.

% ══════════════════════════════════════════════════════════════
\section{Stratificazione e incapsulamento}
% ══════════════════════════════════════════════════════════════

\begin{esercizio}{Quanto pesano gli header}
Un'applicazione invia messaggi di $M$ byte. Lungo la pila TCP/IP su Ethernet si aggiungono un header TCP da 20 byte, un header
IP da 20 byte e 18 byte di header e trailer Ethernet.
\begin{enumerate}[label=(\alph*)]
    \item Che frazione dei byte trasmessi è occupata dagli header se $M = 100$ byte? E se $M = 1460$ byte?
    \item In una sessione remota ogni tasto premuto invia 1 byte. Che frazione è overhead?
    \item Ricavare la formula per una pila di $n$ strati, ognuno dei quali aggiunge $h$ byte.
\end{enumerate}
\end{esercizio}

\svolgimento
Gli header valgono in tutto $h_{\text{tot}} = 20 + 20 + 18 = 58$ byte.

\begin{center}
\begin{tikzpicture}[every node/.style={font=\scriptsize\sffamily}]
  \node[field, fill=llink, minimum width=1.3cm] (e) at (0,0) {Ethernet 14};
  \node[field, fill=lnet, minimum width=1.2cm, right=0pt of e] (i) {IP 20};
  \node[field, fill=ltra, minimum width=1.2cm, right=0pt of i] (t) {TCP 20};
  \node[field, fill=lapp, minimum width=4.2cm, right=0pt of t] (m) {messaggio: $M$ byte};
  \node[field, fill=llink, minimum width=1.2cm, right=0pt of m] {CRC 4};
\end{tikzpicture}
\end{center}

\textbf{(a)}
$$M = 100: \quad \frac{58}{100 + 58} \approx 36{,}7\% \qquad\qquad M = 1460: \quad \frac{58}{1460 + 58} \approx 3{,}8\%$$

\textbf{(b)} $\dfrac{58}{1 + 58} \approx 98{,}3\%$: per spedire un byte se ne trasmettono 59. I messaggi piccoli sono i più
penalizzati dalla stratificazione.

\textbf{(c)} La frazione occupata dagli header è
$$\frac{n\,h}{M + n\,h}$$
che tende a 0 per messaggi grandi e a 1 per messaggi piccoli.

\begin{esercizio}{A quale livello?}
Per ciascuna voce indicare il livello della pila a cinque strati coinvolto e, dove ha senso, l'unità di dati.
\begin{enumerate}
    \item Scegliere il cammino di un pacchetto tra i router.
    \item Consegnare i dati al processo giusto, individuato da un numero di porta.
    \item Delimitare i frame e rilevare gli errori su un singolo collegamento.
    \item Trasformare i bit in segnali elettrici o luminosi.
    \item Chiedere una pagina Web a un server.
    \item Ritrasmettere dati persi lungo il percorso tra i due estremi.
    \item Tradurre un nome come \texttt{www.unina.it} in un indirizzo IP.
    \item Uno switch; un router; un hub (ripetitore).
\end{enumerate}
\end{esercizio}

\svolgimento
\begin{table}[H]
\centering
\begin{tabular}{clp{6.6cm}}
\toprule
\textbf{\#} & \textbf{Livello} & \textbf{Nota} \\
\midrule
1 & Rete & routing; unità: datagramma \\
2 & Trasporto & consegna processo-a-processo; unità: segmento \\
3 & Collegamento & consegna \emph{hop-to-hop}; unità: frame \\
4 & Fisico & unità: bit \\
5 & Applicazione & HTTP; unità: messaggio \\
6 & Trasporto & il trasferimento affidabile è di TCP (end-to-end) \\
7 & Applicazione & il DNS è un'applicazione che offre un servizio alle altre applicazioni \\
8 & Collegamento; Rete; Fisico & lo switch implementa 2 livelli, il router 3, l'hub solo il fisico \\
\bottomrule
\end{tabular}
\caption{Livelli e unità di dati.}
\end{table}

\begin{esercizio}{Che servizio serve?}
Per ciascuna applicazione scegliere il tipo di servizio più adatto tra i sei visti nel capitolo sulla stratificazione, e dire
se ci si aspetta che usi TCP o UDP.
(a) scaricare un file ISO; (b) una telefonata VoIP; (c) una query DNS; (d) stampare un libro diviso in pagine;
(e) pubblicità inviata in massa; (f) un messaggio con conferma di lettura.
\end{esercizio}

\svolgimento
\begin{itemize}
    \item \textbf{(a)} Flusso di byte affidabile (TCP): non può mancare nemmeno un byte.
    \item \textbf{(b)} Connessione inaffidabile (tipicamente UDP): un pacchetto perso è un piccolo ``click'', uno in ritardo è inutile.
    \item \textbf{(c)} Richiesta-risposta (UDP): una domanda, una risposta; se si perde, si ripete.
    \item \textbf{(d)} Flusso di messaggi affidabile: i confini tra le pagine vanno preservati.
    \item \textbf{(e)} Datagram inaffidabile: nessuno si preoccupa se qualche messaggio va perso.
    \item \textbf{(f)} Datagram con conferma: niente connessione, ma il mittente vuole sapere che il messaggio è arrivato.
\end{itemize}

\begin{esercizio}{Chi legge quale header}
L'host A invia una richiesta HTTP al server B. Il percorso attraversa uno switch, poi un router, poi arriva a B. Per ciascun
dispositivo dire fino a quale livello ``apre'' il pacchetto e quali campi legge o modifica.
\end{esercizio}

\svolgimento
\begin{figure}[H]
\centering
\begin{tikzpicture}[every node/.style={font=\scriptsize\sffamily}, lay/.style={minimum width=1.6cm, minimum height=0.45cm, font=\tiny\sffamily}]
  % host A
  \node[lay app, lay] (a5) at (0,0) {Applicazione};
  \node[lay tra, lay, below=0pt of a5] (a4) {Trasporto};
  \node[lay net, lay, below=0pt of a4] (a3) {Rete};
  \node[lay link, lay, below=0pt of a3] (a2) {Collegamento};
  \node[lay phy, lay, below=0pt of a2] (a1) {Fisico};
  \node[lay link, lay] (s2) at (3.3,-1.35) {Collegamento};
  \node[lay phy, lay, below=0pt of s2] (s1) {Fisico};
  \node[lay net, lay] (r3) at (6.6,-0.9) {Rete};
  \node[lay link, lay, below=0pt of r3] (r2) {Collegamento};
  \node[lay phy, lay, below=0pt of r2] (r1) {Fisico};
  \node[lay app, lay] (b5) at (9.9,0) {Applicazione};
  \node[lay tra, lay, below=0pt of b5] (b4) {Trasporto};
  \node[lay net, lay, below=0pt of b4] (b3) {Rete};
  \node[lay link, lay, below=0pt of b3] (b2) {Collegamento};
  \node[lay phy, lay, below=0pt of b2] (b1) {Fisico};
  \foreach \x/\n in {0/host A, 3.3/switch, 6.6/router, 9.9/server B} \node[font=\footnotesize\bfseries\sffamily] at (\x,0.55) {\n};
  \draw[very thick, netred, -{Stealth}, rounded corners=2pt] (0.95,0) -- (0.95,-2.25) -- (2.35,-2.25) -- (2.35,-1.0) -- (4.25,-1.0)
     -- (4.25,-2.25) -- (5.65,-2.25) -- (5.65,-0.55) -- (7.55,-0.55) -- (7.55,-2.25) -- (8.85,-2.25) -- (8.85,0);
\end{tikzpicture}
\caption{Il pacchetto sale lungo la pila solo fino al livello che il dispositivo implementa.}
\end{figure}

\begin{itemize}
    \item \textbf{Switch} (livello 2): legge l'header Ethernet, in particolare il \textbf{MAC di destinazione}, per scegliere la
          porta d'uscita. Non modifica il frame e non guarda dentro il payload.
    \item \textbf{Router} (livello 3): rimuove l'header Ethernet, legge l'\textbf{IP di destinazione} per scegliere l'uscita,
          \textbf{decrementa il TTL} e ricalcola il checksum IP, poi incapsula il datagramma in un \textbf{nuovo frame} con nuovi MAC.
          Non legge l'header TCP né la richiesta HTTP.
    \item \textbf{Server B}: decapsula tutti gli strati fino alla richiesta HTTP, che consegna al processo server individuato
          dalla porta di destinazione (80).
\end{itemize}
Ogni strato tratta il proprio payload come \textbf{opaco}: il router trasporta la richiesta HTTP senza sapere che esiste.

\begin{esercizio}{Simplex, half-duplex o full-duplex?}
Classificare: (a) una trasmissione televisiva; (b) una coppia di walkie-talkie; (c) una telefonata; (d) una tastiera
collegata a un computer; (e) un segmento Ethernet condiviso su hub; (f) un collegamento Ethernet commutato moderno.
\end{esercizio}

\svolgimento
(a) simplex; (b) half-duplex (si parla a turno, ``passo''); (c) full-duplex; (d) simplex; (e) half-duplex (una stazione
alla volta, altrimenti collisione); (f) full-duplex (coppie di fili separate per i due versi, niente collisioni).

\gbox{\iinfo\ Formule da ricordare}{primary!80!black}{
\begin{itemize}
    \item $d_{\text{trasm}} = L/R$; \quad $d_{\text{prop}} = d/s$; \quad lunghezza di un bit sul mezzo $= s/R$.
    \item $P$ pacchetti su $N$ collegamenti store-and-forward: $T = (P + N - 1)\cdot L/R$.
    \item Throughput end-to-end $= \min_i R_i$ (tenendo conto della condivisione).
    \item Maglia completa: $n(n-1)/2$ collegamenti; stella e anello: $n$.
    \item Overhead degli header: $n\,h / (M + n\,h)$.
\end{itemize}
}
